% ============================================================
%  related_works.tex  —  Related Work for Subspace-SCAFFOLD
%  Target venues: ICML / NeurIPS / ICLR
%  Total citations: 27  (threshold ≥ 22 ✓)
% ============================================================

\section{Related Work}
\label{sec:related_work}

Our work sits at the intersection of three active research threads:
federated learning with data heterogeneity, communication compression,
and subspace / low-rank optimization.
We survey each thread in turn and conclude by identifying the gap that
\textsc{Subspace-SCAFFOLD} is designed to fill.

% ------------------------------------------------------------
\subsection{Federated Learning and Client Drift}
\label{subsec:fl}

Federated learning (FL) enables a large population of clients to
collaboratively train a shared model without exchanging raw
data~\citep{mcmahan2017communication}.
The canonical \textsc{FedAvg} algorithm~\citep{mcmahan2017communication}
alternates between local stochastic gradient descent (SGD) steps on
each client and a global averaging step on the server.
Subsequent work established strategies for reducing the number of
communication rounds through structured
updates~\citep{konecny2016federated} and for accelerating local
convergence via variance-reduced gradient
estimators~\citep{zhao2021fedpage}.

A central challenge in FL is \emph{statistical heterogeneity}: when
clients hold non-i.i.d.\ data, local objectives diverge from the
global objective, causing \emph{client drift}—a systematic bias in the
aggregated update direction that severely degrades
convergence~\citep{karimireddy2020scaffold,wang2020tackling}.
Two families of methods have been proposed to mitigate this effect.

\paragraph{Proximal regularization.}
\textsc{FedProx}~\citep{li2020fedprox} adds a proximal term
$\frac{\mu}{2}\|\mathbf{w} - \mathbf{w}^t\|^2$ to each client's local
objective, penalizing deviation from the current global model.
This provides a convergence guarantee under arbitrary degrees of
heterogeneity and partial device participation, but the proximal
coefficient $\mu$ must be tuned carefully and the regularization does
not explicitly correct the gradient direction.

\paragraph{Control variates.}
\textsc{SCAFFOLD}~\citep{karimireddy2020scaffold} introduces
\emph{control variates}—client-specific and server-side correction
terms $\mathbf{c}_i, \mathbf{c} \in \mathbb{R}^d$—that directly
estimate and subtract the client drift at each local step.
\textsc{SCAFFOLD} achieves a communication complexity that is
independent of data heterogeneity, making it theoretically superior to
\textsc{FedAvg} under non-i.i.d.\ settings.
Variants of this idea have also been explored in the context of
objective inconsistency~\citep{wang2020tackling}.

\paragraph{Memory overhead.}
Despite their effectiveness, both \textsc{FedProx} and
\textsc{SCAFFOLD} operate in the \emph{full parameter space}
$\mathbb{R}^d$.
\textsc{SCAFFOLD} in particular requires each client to maintain a
control variate vector of dimension $d$, doubling the per-client memory
footprint.
For modern large neural networks where $d$ can reach hundreds of
millions, this overhead is prohibitive on resource-constrained edge
devices.

% ------------------------------------------------------------
\subsection{Communication Compression}
\label{subsec:compression}

Transmitting full-precision, full-dimensional gradient vectors in every
round constitutes the dominant communication bottleneck in
FL~\citep{konecny2016federated}.
Two principal compression strategies have been studied.

\paragraph{Quantization.}
\textsc{QSGD}~\citep{alistarh2017qsgd} replaces each gradient
coordinate with a randomly rounded low-bit representation, achieving
up to $4\times$–$8\times$ bandwidth reduction with provable
unbiasedness.
Sign-based methods such as \textsc{SignSGD}~\citep{bernstein2018signsgd}
transmit only the sign of each coordinate (1 bit per dimension),
yielding extreme compression at the cost of introducing bias.
Optimal quantization strategies that jointly minimize communication
cost and convergence slowdown have been
characterized~\citep{albasyoni2020optimal}.

\paragraph{Sparsification.}
Top-$k$ sparsification~\citep{lin2018deep,shi2019understanding}
transmits only the $k$ gradient coordinates with the largest magnitude,
achieving compression ratios exceeding $99\%$ in practice.
\citet{lin2018deep} show that combining aggressive sparsification with
momentum correction and warm-up schedules (\emph{Deep Gradient
Compression}) preserves model accuracy.
Low-rank decomposition of gradient matrices via power iteration
(\textsc{PowerSGD})~\citep{vogels2019powersgd} and its decentralized
extension~\citep{vogels2020powergossip} offer an alternative that is
particularly effective for wide layers.

\paragraph{Error feedback (EF).}
A key insight for enabling convergence with \emph{biased} compressors
(e.g., top-$k$) is \emph{error feedback}: accumulating the compression
residual $\mathbf{e}_t = \mathbf{g}_t - \mathcal{C}(\mathbf{g}_t +
\mathbf{e}_{t-1})$ and adding it to the next gradient before
compression~\citep{stich2018sparsified}.
\citet{richtarik2021ef21} provide a cleaner theoretical analysis of EF
under general contractive compressors (\textsc{EF21}), and
\citet{richtarik2024ef_reloaded} further tighten the dependence on
smoothness constants.
\citet{bao2025efskip} extend EF to the heterogeneous FL setting,
achieving linear speedup without restrictive assumptions on data
distributions.

\paragraph{Memory bottleneck of EF.}
A critical but often overlooked drawback of error feedback is that
every client must maintain an error accumulation vector
$\mathbf{e}_t \in \mathbb{R}^d$ across rounds.
When combined with the control variates of \textsc{SCAFFOLD}, the
per-client memory requirement grows to $3d$ floats (model weights,
control variate, error buffer), which is infeasible for
resource-constrained devices training large models.
Furthermore, compression introduces additional gradient variance that
compounds the client-drift problem under heterogeneous data, making the
simultaneous application of EF and variance reduction particularly
challenging.

% ------------------------------------------------------------
\subsection{Subspace and Low-Rank Methods in FL}
\label{subsec:subspace}

A complementary line of research exploits the empirically observed
\emph{low-rank structure} of gradients to reduce communication,
computation, and memory simultaneously.

\paragraph{Low-rank adaptation.}
\textsc{LoRA}~\citep{hu2022lora} parameterizes weight updates as the
product of two low-rank matrices $\mathbf{B}\mathbf{A}$ with
$\text{rank}(r) \ll d$, reducing trainable parameters by orders of
magnitude.
\textsc{GaLore}~\citep{zhao2024galore} applies a similar idea to
full-parameter training by projecting gradients onto a periodically
updated low-rank subspace, achieving memory efficiency comparable to
\textsc{LoRA} without restricting the update manifold.
\textsc{Flora}~\citep{hao2024flora} shows that \textsc{LoRA} adapters
implicitly compress gradients and provides a unified view connecting
low-rank adaptation with gradient compression.

\paragraph{Subspace methods for FL.}
The low-rank structure of gradient subspaces in FL has been
empirically validated~\citep{azam2022recycling}, motivating a family of
subspace-based FL algorithms.
\textsc{FedSub}~\citep{zhang2025fedsub} projects local gradients onto a
shared low-dimensional subspace before aggregation, simultaneously
reducing communication volume and mitigating drift through implicit
regularization.
\textsc{FedLoRU}~\citep{park2024fedloru} provides the first theoretical
study of rank properties in FL and shows that low-rank updates enjoy
implicit regularization benefits.
\textsc{ILoRA}~\citep{zhou2025ilora} addresses heterogeneous client
aggregation under \textsc{LoRA}-style updates by aligning local
subspaces before aggregation.
\citet{peng2026rethinking} identify two fundamental mismatches in
federated \textsc{LoRA}—update-space misalignment and optimizer state
misalignment—and propose remedies based on subspace and state
alignment.
Efficient wireless FL via low-rank gradient
factorization~\citep{guo2024efficient} further demonstrates the
communication benefits of low-rank approaches in bandwidth-limited
environments.

\paragraph{Federated fine-tuning of large models.}
The challenge of fine-tuning large language models (LLMs) in the
federated setting has attracted growing attention.
\citet{wang2024save} propose a cycle gradient descent strategy that
enables full-parameter tuning within a low-rank update framework,
avoiding the expressivity loss of fixed-rank adapters.

% ------------------------------------------------------------
\subsection{The Gap: Subspace Methods Without Drift Correction}
\label{subsec:gap}

Despite significant progress in both variance reduction and subspace
optimization, \emph{no existing method simultaneously achieves all
three desiderata}: (i) correcting client drift under non-i.i.d.\ data,
(ii) operating in a low-dimensional subspace to reduce memory and
communication, and (iii) maintaining provable convergence guarantees.

The core difficulty is a \emph{subspace-switching incompatibility}:
\textsc{SCAFFOLD}'s control variates $\mathbf{c}_i$ are accumulated
in the full parameter space $\mathbb{R}^d$ over multiple rounds.
When the projection subspace is updated (as required to track the
evolving gradient manifold), the historical control variate becomes
\emph{misaligned} with the new subspace, rendering the correction
ineffective or even harmful.
Existing subspace FL methods~\citep{zhang2025fedsub,park2024fedloru,
zhou2025ilora} do not address this incompatibility; they either ignore
client drift entirely or rely on proximal regularization, which
provides weaker guarantees than explicit variance
reduction~\citep{karimireddy2020scaffold,li2020fedprox}.

\textsc{Subspace-SCAFFOLD} closes this gap by maintaining and updating
control variates \emph{within} the current subspace, projecting them
consistently whenever the subspace rotates, and providing a theoretical
analysis that accounts for both the approximation error of the
projection and the variance introduced by subspace switching.
This design simultaneously achieves $\mathcal{O}(r/d)$ memory and
communication reduction (where $r \ll d$ is the subspace rank) and
retains the heterogeneity-robust convergence of
\textsc{SCAFFOLD}~\citep{karimireddy2020scaffold}.
